\capitulocapa{images/cap-10.png}{Como lembrar preferências do usuário?}
\label{chap:persistencia}

\section{A configuração que durava até fechar a aba}

\begin{historia}
A tela de Configurações do DevPanel ganhou duas preferências: tema claro ou
escuro e densidade confortável ou compacta. A equipe colocou ambas numa store e
o resultado pareceu pronto. Header, Sidebar e tabelas reagiam imediatamente.

Na manhã seguinte, uma pessoa abriu o painel e encontrou tudo no padrão. Escolheu
novamente o tema escuro e a densidade compacta. Depois de atualizar a página
durante um suporte, perdeu as escolhas pela terceira vez.

O pedido chegou como uma frase curta: ``O painel precisa lembrar de mim.'' Na
mesma store também vivia o estado de abertura do Drawer de notificações.
Persistir a store inteira lembraria mais do que a pessoa pediu: o Drawer poderia
reaparecer aberto após cada visita.

A equipe percebeu que a decisão não era apenas como salvar estado. Era escolher
qual estado merecia sobreviver.
\end{historia}

\imagemcapitulo
  {ch10-fluxo-persist-review}
  {O middleware persiste preferências e exclui o estado transitório do Drawer.}
  {fluxo-persist}

Uma store comum vive na memória do processo da página. Atualizar ou fechar a
página encerra essa memória. Persistência adiciona outro ciclo de vida: parte do
estado pode ser escrita num armazenamento e recuperada numa execução futura.

\section{Memória e persistência resolvem tempos diferentes}

Enquanto o DevPanel está aberto, Zustand mantém a fotografia atual e notifica os
componentes interessados. Isso é memória de execução. Ela é rápida e suficiente
para modais, Drawers, filtros temporários e passos de um fluxo.

Preferências possuem outra expectativa. Tema e densidade representam escolhas
que a pessoa não quer repetir a cada carregamento. O navegador pode guardá-las
entre sessões, mas a store continua sendo a interface usada pela aplicação.

Persistir não transforma o storage na nova store. Os componentes não devem ler
\codigo{localStorage} diretamente. A store permanece como fonte de verdade em
tempo de execução; o armazenamento fornece uma fotografia que pode restaurá-la.

\begin{decisao}
Persistência é uma decisão de ciclo de vida. Antes de escolher uma API, pergunte
se a informação deve sobreviver à página, à aba, ao dispositivo ou à conta.
\end{decisao}

\section{A store antes da persistência}

O contrato reúne as duas preferências e um estado transitório da interface:

\begin{lstlisting}[caption={Preferências e Drawer possuem ciclos de vida diferentes.}]
type Theme = 'light' | 'dark'
type Density = 'comfortable' | 'compact'

type PreferenceStore = {
  theme: Theme
  density: Density
  isNotificationDrawerOpen: boolean
  selectTheme: (theme: Theme) => void
  selectDensity: (density: Density) => void
  openNotificationDrawer: () => void
  closeNotificationDrawer: () => void
}
\end{lstlisting}

Os três valores são estado do cliente e podem viver na mesma store neste estágio
do projeto. Isso não significa que devam ter a mesma duração. Tema e densidade
são preferências; Drawer aberto descreve um instante da interface.

Sem persistência, a implementação é conhecida:

\begin{lstlisting}
export const usePreferenceStore =
  create<PreferenceStore>()((set) => ({
    theme: 'light',
    density: 'comfortable',
    isNotificationDrawerOpen: false,
    selectTheme: (theme) => set({ theme }),
    selectDensity: (density) => set({ density }),
    openNotificationDrawer: () =>
      set({ isNotificationDrawerOpen: true }),
    closeNotificationDrawer: () =>
      set({ isNotificationDrawerOpen: false }),
  }))
\end{lstlisting}

A store funciona, mas sempre recomeça com os valores declarados na criadora.

\section{Adicionando o middleware persist}

Um middleware envolve a função criadora e acrescenta um comportamento ao ciclo
da store. Para persistência, importamos \codigo{persist}:

\begin{lstlisting}[caption={Store de preferências persistida.},label={lst:preference-persist}]
import { create } from 'zustand'
import { persist } from 'zustand/middleware'

export const usePreferenceStore =
  create<PreferenceStore>()(
    persist(
      (set) => ({
        theme: 'light',
        density: 'comfortable',
        isNotificationDrawerOpen: false,

        selectTheme: (theme) => set({ theme }),
        selectDensity: (density) => set({ density }),
        openNotificationDrawer: () =>
          set({ isNotificationDrawerOpen: true }),
        closeNotificationDrawer: () =>
          set({ isNotificationDrawerOpen: false }),
      }),
      {
        name: 'devpanel-preferences',
      },
    ),
  )
\end{lstlisting}

\codigo{name} identifica a entrada no armazenamento e deve ser único para a
aplicação. No navegador, o middleware usa \codigo{localStorage} por padrão e
serializa os dados em JSON.

Quando uma ação muda o estado, \codigo{persist} grava a fotografia correspondente.
Quando a store é criada novamente, o middleware procura essa chave e recupera o
conteúdo disponível.

Essa primeira versão ainda contém um problema: sem outra configuração, estamos
dispostos a persistir dados demais.

\section{Persistir somente o que foi escolhido}

\codigo{partialize} recebe o estado atual e devolve a parte que será armazenada:

\begin{lstlisting}[caption={Lista explícita das preferências persistidas.},label={lst:preference-partialize}]
{
  name: 'devpanel-preferences',
  partialize: (state) => ({
    theme: state.theme,
    density: state.density,
  }),
}
\end{lstlisting}

Agora o storage contém tema e densidade. Ações não precisam ser serializadas e
\codigo{isNotificationDrawerOpen} fica de fora. Ao recarregar a página, o Drawer
usa o valor inicial \codigo{false}, mesmo que estivesse aberto antes.

Preferimos declarar o que entra em vez de listar o que deve sair. Quando um novo
campo for adicionado à store, ele não ganhará vida longa por acidente. A equipe
precisará decidir conscientemente se esse campo pertence ao contrato persistido.

\imagemcapitulo
  {ch10-quadro-dados-persistidos}
  {Cada dado exige duração e destino compatíveis com quem realmente o possui.}
  {quadro-dados-persistidos}

\section{Storage não é banco de dados nem cofre}

\codigo{localStorage} pertence à origem do site e permanece após fechar o
navegador. Ele combina com preferências pequenas que não são sigilosas. Também
é possível escolher outro mecanismo, por exemplo \codigo{sessionStorage}:

\begin{lstlisting}
import {
  createJSONStorage,
  persist,
} from 'zustand/middleware'

persist(createPreferenceState, {
  name: 'devpanel-preferences',
  storage: createJSONStorage(() => sessionStorage),
})
\end{lstlisting}

Nesse caso, a informação dura durante a sessão da aba. Armazenamentos assíncronos
também podem ser integrados, mas mudam o momento em que a restauração termina.

Persistência local não sincroniza preferências entre computadores e não substitui
uma API quando a escolha pertence à conta. Também não torna o conteúdo seguro.
Scripts executados na mesma origem podem acessar o storage. Tokens e credenciais
exigem uma decisão de segurança própria, que será discutida no Capítulo~16.

\section{O formato salvo também evolui}

Meses depois, a equipe pode renomear campos, trocar valores aceitos ou reorganizar
a store. Pessoas que já usaram o DevPanel continuarão com a fotografia antiga no
navegador.

Imagine que a primeira versão persistia \codigo{compactMode} como booleano. A
versão atual usa \codigo{density} com valores semânticos. Apenas alterar o tipo
no código não converte o JSON já gravado.

O contrato persistido precisa de uma versão:

\begin{lstlisting}[caption={Migração do formato antigo para a versão atual.},label={lst:preference-migrate}]
type PersistedVersionZero = {
  theme: Theme
  compactMode: boolean
}

{
  name: 'devpanel-preferences',
  version: 1,
  partialize: (state) => ({
    theme: state.theme,
    density: state.density,
  }),
  migrate: (persistedState, version) => {
    if (version === 0) {
      const oldState =
        persistedState as PersistedVersionZero

      return {
        theme: oldState.theme,
        density: oldState.compactMode
          ? 'compact'
          : 'comfortable',
      }
    }

    return persistedState
  },
}
\end{lstlisting}

\codigo{version: 1} descreve o formato atual. Quando a versão armazenada é
diferente, \codigo{migrate} recebe o conteúdo antigo e sua versão. A função deve
devolver dados compatíveis com o contrato atual.

Uma migração não é necessária para qualquer alteração interna da store. Ela é
necessária quando muda o formato da parcela persistida. Adicionar uma ação, por
exemplo, não altera o JSON selecionado por \codigo{partialize}.

Migrações devem ser pequenas, determinísticas e capazes de lidar com dados
antigos. Se uma fotografia não puder ser convertida com segurança, restaurar os
valores padrão pode ser melhor que iniciar a aplicação num estado inválido.

\section{Hidratação: o encontro de duas fotografias}

Hidratação é o processo de ler o estado persistido e combiná-lo com o estado
inicial da store. No DevPanel executado com Vite e \codigo{localStorage}, o
armazenamento é síncrono e a restauração acontece durante a criação da store.

O fluxo possui duas fontes temporárias:

\begin{enumerate}
  \item o código cria tema claro, densidade confortável e Drawer fechado;
  \item \codigo{persist} procura a chave no storage;
  \item encontra tema escuro e densidade compacta da sessão anterior;
  \item combina os dados persistidos com o estado atual;
  \item ações e campos não persistidos continuam vindo da criadora.
\end{enumerate}

A combinação padrão é rasa. Se persistirmos parcialmente um objeto aninhado, a
parcela antiga pode substituir o objeto atual inteiro. Por isso, preferências
planas tornam esse contrato mais previsível; estruturas aninhadas podem exigir
uma função de combinação explícita.

Com um armazenamento assíncrono, a primeira leitura pode ocorrer antes do fim da
hidratação. Interfaces que dependem da preferência no primeiro quadro precisam
considerar um estado de ``ainda não hidratado''. Em aplicações renderizadas no
servidor, o assunto também envolve consistência entre HTML inicial e cliente;
vamos tratar essa fronteira no Capítulo~17.

\section{O que não deve ser persistido}

Uma lista curta ajuda o time durante a revisão:

\begin{itemize}
  \item não persista valores que podem ser derivados de outros já salvos;
  \item não persista respostas remotas apenas para improvisar um cache;
  \item não persista estados transitórios sem expectativa clara de retomada;
  \item não persista dados sensíveis por conveniência;
  \item não persista campos ``talvez úteis'' sem política de expiração e migração.
\end{itemize}

Um rascunho de formulário pode merecer recuperação. Um modal aberto normalmente
não merece. A diferença não está no tipo TypeScript; ambos podem ser booleanos
ou strings. Está na experiência que o produto promete quando a pessoa retorna.

Persistência também cria responsabilidade operacional. Dados antigos permanecem
em navegadores que a equipe não controla. Cada campo salvo aumenta o contrato de
compatibilidade das próximas versões.

\section{Exercício: lembre a densidade, esqueça o Drawer}

\begin{tarefa}
O produto pediu que a densidade compacta continue selecionada depois de fechar e
abrir o navegador. O Drawer de notificações, porém, deve sempre iniciar fechado.
Adicione \codigo{persist} à store, use \codigo{partialize} para declarar os dados
permitidos e confirme o comportamento após recarregar a página.
\end{tarefa}

Depois da implementação, registre:

\begin{enumerate}
  \item qual chave foi criada no storage;
  \item quais campos aparecem no JSON persistido;
  \item por que o Drawer volta fechado;
  \item qual mudança futura exigiria uma migração;
  \item qual preferência deveria viver na conta, e não apenas no navegador.
\end{enumerate}

\espacoresposta{9}

\begin{resumo}
\begin{itemize}
  \item Estado em memória e estado persistido possuem ciclos de vida distintos.
  \item \codigo{persist} conecta a store a um mecanismo de armazenamento.
  \item \codigo{partialize} declara a parcela autorizada a sobreviver.
  \item Storage local é adequado para preferências pequenas, não para segredos.
  \item \codigo{version} e \codigo{migrate} protegem mudanças de formato.
  \item Hidratação combina a fotografia persistida com o estado atual.
  \item Persistir um campo cria um contrato de experiência e compatibilidade.
\end{itemize}
\end{resumo}

Com preferências persistidas, o DevPanel agora possui mudanças que atravessam
sessões. Na próxima parte, começaremos a torná-las observáveis e a organizar as
stores para que o time consiga manter esse crescimento.
